\begin{figure}[tbp]
\centering
\begin{minipage}[t]{.48\linewidth}
\centering\ViewHeading{(a) Exact dimension formula for $d=2$}
\begin{tikzpicture}
\begin{axis}[viewaxis,width=7cm,height=5cm,xmin=0,xmax=4,ymin=0,ymax=2.25,
 xlabel={$\alpha$},ylabel={$h(\alpha)$},xtick={.5,1,2,3,4},ytick={0,1,2}]
\path[fill=viewgray!10] (axis cs:0,0) rectangle (axis cs:.5,2.25);
\addplot[viewblue,domain=.501:4,samples=70] {3/(1+x)};
\addplot[only marks,mark=o,viewblue] coordinates {(.5,2)};
\draw[dashed,viewgray] (axis cs:.5,0)--(axis cs:.5,2.25);
\node[font=\scriptsize,rotate=90] at (axis cs:.23,1.1) {outside stated theorem};
\node[font=\footnotesize,viewblue] at (axis cs:2.4,1.8) {$h(\alpha)=3/(1+\alpha)$};
\end{axis}
\end{tikzpicture}
\end{minipage}\hfill
\begin{minipage}[t]{.48\linewidth}
\centering\ViewHeading{(b) Dense and null are compatible}
\begin{tikzpicture}
\node[vbox,text width=5.9cm] at (0,0)
 {$Z=\iota_{\boldsymbol\lambda}(R_{\alpha,C}\setminus\{0\})$\\
 $\overline Z=\mathcal M_{\boldsymbol\lambda}$,\quad $\mathcal H^d(Z)=0$};
\node[vresult,text width=5.9cm] at (0,-1.5)
 {Hausdorff dimension: $h=(d+1)/(1+\alpha)$\\Both box dimensions: $d$};
\node[vinput,text width=5.9cm] at (0,-3.05)
 {Codimension $\eta=d-h$ is not a width.\\Radii $\lambda_j$ and offset $r$ are separate inputs.};
\end{tikzpicture}
\end{minipage}
\medskip
\begin{tabular}{@{}llll@{}}
\toprule
Quantity & $b_q$ & $\Fill_{q,x}(1)$ & $\widehat{\Fill}_{q,x}(1)$\\\midrule
Threshold & $<Cq^{-\alpha}$ & $>C^{-1}q^\alpha$ & $>C^{-1}q^{\alpha+1}$\\\bottomrule
\end{tabular}
\par\smallskip{\small Normalized component: $\|q^{-1}v_q\|_2<Cq^{-(\alpha+1)}$.}
\caption{Dimension, density, and threshold normalization. The curve is
the stated theorem's exact formula, not a fitted numerical estimate.
Dense subsets and their closures have the same finite closed-ball
covering numbers here, while their Hausdorff dimensions differ. No
positive-thickness fractal region is implied.}
\label{fig:arithmetic-dimensions}
\end{figure}